\makeatletter
\@ifundefined{keepfig}{\newenvironment{keepfig}{\par\edef\keepfig@pi{\the\parindent}\noindent
  \begin{minipage}{\linewidth}\parindent=\keepfig@pi\relax}{\end{minipage}\par}}{}
\makeatother
\section{Группы автоморфизмов полей}

Пусть $E\supset F$~--- поля. $F$-изоморфизм $E\to E$ называется \emph{$F$-автоморфизмом} поля $E$. Композиция
$F$-автоморфизмов~--- снова $F$-автоморфизм, и обратное отображение к $F$-автоморфизму~--- тоже $F$-автоморфизм (обратное
к биективному гомоморфизму колец является гомоморфизмом и тождественно на $F$). Поэтому $F$-автоморфизмы поля $E$
образуют группу относительно композиции; она обозначается $\Aut(E/F)$. Через $\Aut(E)$ обозначим группу \emph{всех}
автоморфизмов поля $E$. Каждый автоморфизм тождествен на простом подполе $F_0\subset E$ (он переводит $1$ в $1$, а значит,
и суммы единиц, их разности и частные), так что $\Aut(E)=\Aut(E/F_0)$.

Главные вопросы этого раздела: насколько велика группа $\Aut(E/F)$ по сравнению со степенью $[E:F]$ и можно ли по группе
автоморфизмов восстановить поле $F$.

\begin{example}[3.1]
(a) У поля $\CC$ есть два очевидных автоморфизма: тождественный и комплексное сопряжение. Других
$\R$-автоморфизмов нет: $\CC=\R[i]$, и $\R$-гомоморфизм $\CC\to\CC$ определяется образом $i$, который должен быть
корнем $X^2+1$, т.\,е.\ равен $\pm i$ (\rf{2.1}(b)). Итак, $\Aut(\CC/\R)=\{\id,\ z\mapsto\bar z\}$. С помощью аксиомы
выбора можно построить несчётно много автоморфизмов $\CC$ (\rf{9.18}); все они, кроме этих двух, разрывны:
непрерывный автоморфизм тождествен на $\Q$, а значит, по непрерывности и на $\R$, т.\,е.\ лежит в $\Aut(\CC/\R)$.

(b) Пусть $E=\CC(X)$~--- поле рациональных функций. $\CC$-автоморфизм поля $E$ переводит $X$ в другую образующую $E$
над $\CC$, а образующие~--- это в точности дроби $\frac{aX+b}{cX+d}$ с $ad-bc\ne0$ (\rf{3.1b} ниже). Поэтому $\Aut(E/\CC)$
состоит из отображений $f(X)\mapsto f\bigl(\frac{aX+b}{cX+d}\bigr)$, $ad-bc\ne0$, и
\[
\Aut(\CC(X)/\CC)\cong PGL_2(\CC)
\]
--- группе обратимых комплексных матриц $2\times2$ по модулю её центра (скалярных матриц).

(c) Группа $\Aut(\CC(X_1,X_2)/\CC)$ устроена гораздо сложнее (\rf{3.1c}).
\end{example}

Утверждение (b) верно над любым полем $k$ вместо $\CC$. Докажем его; основной шаг~--- вычисление степени расширения
$k(X)/k(u)$.

\begin{lemma}[3.1a]
Пусть $k$~--- поле, $X$ трансцендентен над $k$, и $u\in k(X)\smallsetminus k$. Запишем $u=g/h$, где $g,h\in k[X]$ взаимно
просты, и положим $n=\max(\deg g,\deg h)$. Тогда $u$ трансцендентен над $k$, элемент $X$ алгебраичен над $k(u)$ и
\[
[k(X):k(u)]=n;
\]
минимальный многочлен $X$ над $k(u)$ с точностью до множителя из $k(u)^\times$ равен $P(T)=g(T)-u\,h(T)$.
\end{lemma}

\begin{proof}
Пусть $g=\sum g_iT^i$, $h=\sum h_iT^i$. Многочлен $P(T)=g(T)-uh(T)\in k(u)[T]$ имеет корень $X$: $P(X)=g(X)-\frac{g(X)}{h(X)}h(X)=0$.
Его коэффициент при $T^n$ равен $g_n-uh_n$, где $g_n,h_n$ не равны нулю одновременно. Если бы $g_n-uh_n=0$, то $h_n\ne0$ и
$u=g_n/h_n\in k$~--- противоречие. Значит, $\deg_TP=n$, и $X$ алгебраичен над $k(u)$. Если бы $u$ был алгебраичен над
$k$, то $X$ был бы алгебраичен над $k$ (\rf{1.31}); поэтому $u$ трансцендентен, и $k[u]$~--- кольцо многочленов от $u$,
а $k[u][T]=k[u,T]$~--- кольцо многочленов от двух переменных, факториальное (\rf{1.16c}).

Покажем, что $P$ неприводим в $k(u)[T]$. Пусть $P=AB$, где $A,B\in k(u)[T]$ имеют положительную степень по $T$. По лемме
Гаусса (\rf{1.16c}(a), $R=k[u]$) можно считать, что $A,B\in k[u,T]$. Степень $P$ по $u$ равна $1$ (так как $h\ne0$), а степени
по $u$ при умножении складываются; значит, один из множителей, скажем $B$, не зависит от $u$: $B\in k[T]$, $\deg B\ge1$.
Запишем $A=\sum_jA_j(T)u^j$; сравнивая коэффициенты при степенях $u$ в равенстве $g(T)-u\,h(T)=B(T)\sum_jA_j(T)u^j$,
получаем $g=BA_0$ и $-h=BA_1$, т.\,е.\ $B$~--- общий делитель $g$ и $h$ положительной степени. Это противоречит
взаимной простоте $g$ и $h$.

Итак, $P$ неприводим, и $P/(g_n-uh_n)$~--- минимальный многочлен $X$ над $k(u)$. Так как $k(X)=k(u)(X)=k(u)[X]$, получаем
$[k(X):k(u)]=\deg_TP=n$ (\rf{1.27c}).
\end{proof}

Число $n$ называют \emph{степенью} рациональной функции $u$. Например, $[k(X):k(X^2)]=2$, $[k(X):k(X+1/X)]=2$, а
$[k(X):k(u)]=6$ для $u=\frac{(X^2-X+1)^3}{X^2(X-1)^2}$ (числитель и знаменатель взаимно просты: многочлен $X^2-X+1$
равен $1$ в точках $0$ и $1$).

\begin{proposition}[3.1b]
Пусть $k$~--- поле.
\begin{enumerate}[label=\textup{(\alph*)}]
\item Элемент $u\in k(X)$ порождает $k(X)$ над $k$ (т.\,е.\ $k(u)=k(X)$) тогда и только тогда, когда
$u=\frac{aX+b}{cX+d}$, где $a,b,c,d\in k$ и $ad-bc\ne0$.
\item Для каждого такого $u$ существует единственный $k$-автоморфизм $\sigma$ поля $k(X)$ с $\sigma(X)=u$, а именно
$f(X)\mapsto f(u)$; других $k$-автоморфизмов у $k(X)$ нет.
\item Для $M=\begin{pmatrix}a&b\\c&d\end{pmatrix}\in GL_2(k)$ положим $M\cdot X=\frac{aX+b}{cX+d}$ и
$\sigma_M(f)(X)=f(M\cdot X)$. Тогда $\sigma_M\sigma_N=\sigma_{NM}$, отображение $M\mapsto\sigma_{M^{-1}}$~---
сюръективный гомоморфизм $GL_2(k)\to\Aut(k(X)/k)$, а его ядро состоит из скалярных матриц. Следовательно,
\[
\Aut(k(X)/k)\cong GL_2(k)/k^\times=PGL_2(k).
\]
\end{enumerate}
\end{proposition}

\begin{proof}
(a) Если $k(u)=k(X)$, то $u\notin k$ и по \rf{3.1a} степень $u$ равна $[k(X):k(u)]=1$, т.\,е.\ $u=\frac{aX+b}{cX+d}$ с
взаимно простыми числителем и знаменателем, не являющимися оба константами. Если бы $ad-bc=0$, то строки $(a,b)$ и
$(c,d)$ были бы пропорциональны, т.\,е.\ один из многочленов $aX+b$, $cX+d$ был бы кратен другому, и $u$ был бы константой
(или знаменатель был бы нулём). Обратно, пусть $ad-bc\ne0$. Тогда $u\notin k$: из $aX+b=\lambda(cX+d)$ следовало бы
$a=\lambda c$, $b=\lambda d$ и $ad-bc=0$. Из $u(cX+d)=aX+b$ получаем $X(a-cu)=du-b$, причём $a-cu\ne0$ (иначе $c=a=0$ и
$ad-bc=0$). Значит, $X=\frac{du-b}{a-cu}\in k(u)$, и $k(u)=k(X)$.

(b) По \rf{2.1}(a) $k$-гомоморфизмы $k(X)\to k(X)$ взаимно однозначно соответствуют трансцендентным над $k$ элементам
$u\in k(X)$: гомоморфизм, отвечающий $u$,~--- это $f(X)\mapsto f(u)$, и его образ равен $k(u)$. Он является
автоморфизмом тогда и только тогда, когда $k(u)=k(X)$; остаётся применить (a) (элементы вида (a) трансцендентны по
\rf{3.1a}).

(c) Имеем $\sigma_N(f)=f(N\cdot X)$, и, применяя $\sigma_M$ (подстановку $X\mapsto M\cdot X$), получаем
$\sigma_M\sigma_N(f)=f\bigl(N\cdot(M\cdot X)\bigr)$. Прямое вычисление показывает, что композиция дробно-линейных
подстановок отвечает произведению матриц: $N\cdot(M\cdot X)=(NM)\cdot X$. Значит, $\sigma_M\sigma_N=\sigma_{NM}$ (порядок
меняется, потому что подстановка действует <<изнутри>>), и для $\rho(M)=\sigma_{M^{-1}}$ имеем
$\rho(M)\rho(N)=\sigma_{N^{-1}M^{-1}}=\sigma_{(MN)^{-1}}=\rho(MN)$. По (a) и (b) $\rho$ сюръективен. Наконец,
$\sigma_M=\id$ $\iff$ $\frac{aX+b}{cX+d}=X$ $\iff$ $cX^2+(d-a)X-b=0$ $\iff$ $c=b=0$, $a=d$, т.\,е.\ $M$ скалярна.
\end{proof}

\begin{aside}[3.1c]
(1) \emph{Сфера Римана.} Аналитики заметят, что $PGL_2(\CC)$~--- это и группа (биголоморфных) автоморфизмов сферы
Римана $\mathbb P^1(\CC)=\CC\cup\{\infty\}$: это дробно-линейные отображения $z\mapsto\frac{az+b}{cz+d}$. Объяснение таково:
мероморфные функции на сфере Римана~--- это в точности рациональные функции от $z$, т.\,е.\ их поле есть $\CC(z)\cong\CC(X)$,
а автоморфизм $\varphi$ сферы действует на этом поле по правилу $f\mapsto f\circ\varphi$. Это сопоставление~--- биекция
$\Aut(\mathbb P^1(\CC))\to\Aut(\CC(z)/\CC)$ (переставляющая порядок композиции, как в \rf{3.1b}(c)). Эти факты комплексного
анализа мы не доказываем; алгебраическая часть картины~--- предложение \rf{3.1b}.

(2) \emph{Две и более переменных.} Обратимая матрица $A=(a_{ij})_{0\le i,j\le2}$ задаёт $\CC$-автоморфизм поля
$\CC(X_1,X_2)$:
\[
X_i\ \mapsto\ \frac{a_{i0}+a_{i1}X_1+a_{i2}X_2}{a_{00}+a_{01}X_1+a_{02}X_2}\qquad(i=1,2)
\]
(в однородных координатах $X_i=x_i/x_0$ это линейная замена $x\mapsto Ax$), и так получается вложение
$PGL_3(\CC)=\Aut(\mathbb P^2(\CC))\hookrightarrow\Aut(\CC(X_1,X_2)/\CC)$. Оно далеко от сюръективности. Например,
$\tau\colon X_1\mapsto X_1$, $X_2\mapsto X_2+X_1^2$~--- автоморфизм (он обратим уже на кольце $\CC[X_1,X_2]$: обратный
$X_2\mapsto X_2-X_1^2$), но не дробно-линейный: если $X_2+X_1^2=L_1/L_0$ с многочленами $L_0\ne0$, $L_1$ степени $\le1$, то
$L_0\cdot(X_2+X_1^2)=L_1$, а степень левой части не меньше $2$. Ещё один пример~--- инволюция $X_i\mapsto1/X_i$. Группа
$\Aut(\CC(X_1,\dots,X_n)/\CC)$ называется \emph{группой Кремоны} (при $n=1$ это группа $PGL_2(\CC)$ из \rf{3.1b}); это группа бирациональных автоморфизмов
проективного пространства $\mathbb P^n(\CC)$, её изучение~--- часть алгебраической геометрии. Уже при $n=2$ она очень
велика, а при $n\ge3$ её строение известно гораздо хуже.
\end{aside}

\begin{example}[3.1d]
$\Aut(\R)=\{\id\}$. Действительно, пусть $\sigma\in\Aut(\R)$. Число $x\ge0$ является квадратом, $x=y^2$, поэтому
$\sigma(x)=\sigma(y)^2\ge0$; значит, если $x\le x'$, то $\sigma(x')-\sigma(x)=\sigma(x'-x)\ge0$, т.\,е.\ $\sigma$ монотонен. Кроме того, $\sigma$
тождествен на $\Q$. Если $r<x<s$ с рациональными $r,s$, то $r<\sigma(x)<s$; так как $r$ и $s$ можно взять сколь угодно
близкими к $x$, получаем $\sigma(x)=x$. Отсюда, в частности, каждый автоморфизм $\CC$, переводящий $\R$ в себя, лежит в
$\Aut(\CC/\R)$, т.\,е.\ это тождество или сопряжение.
\end{example}

Далее нас будут интересовать группы $\Aut(E/F)$ для \emph{конечных} расширений $E/F$.

\begin{proposition}[3.2]
Пусть $E$~--- поле разложения сепарабельного многочлена $f\in F[X]$. Тогда $|\Aut(E/F)|=[E:F]$.
\end{proposition}

\begin{proof}
Так как $f$ сепарабелен, он имеет в $E$ ровно $\deg f$ различных корней, т.\,е.\ раскладывается в $E[X]$ на попарно
различные линейные множители. Поле $E$ порождено корнями $f$ и расщепляет $f$, поэтому по \rf{2.12} (с $\Om=E$) число
$F$-гомоморфизмов $E\to E$ равно $[E:F]$. Поле $E$ конечно над $F$ (\rf{1.30}), а каждый $F$-гомоморфизм $E\to E$~---
инъективное $F$-линейное отображение конечномерного пространства в себя, т.\,е.\ биекция. Значит, все эти гомоморфизмы~---
автоморфизмы.
\end{proof}

\begin{keepfig}
\begin{center}
\begin{tikzcd}[column sep=huge,row sep=large]
E\arrow[r,cnew,dashed,"\varphi"] & E\\
F\arrow[u,hook]\arrow[ur,cold,hook] &
\end{tikzcd}
\end{center}
\figcap{К доказательству \rf{3.2}: левое $E$ порождено корнями $f$, правое расщепляет $f$ на различные линейные множители.
По \rf{2.12} таких $\varphi$ ровно $[E:F]$, и каждый из них~--- автоморфизм, так как $E$ конечномерно над $F$.}
\end{keepfig}

Обе части условия существенны: <<поле разложения>> обеспечивает, что всем корням есть куда переходить, а
<<сепарабельный>>~--- что корней достаточно много.

\begin{example}[3.3]
Пусть $E=F[\alpha]$~--- простое алгебраическое расширение, и $f\in F[X]$~--- многочлен с корнем $\alpha$. Каждый
$F$-автоморфизм $E$ переводит $\alpha$ в корень $f$ (\rf{2.1}(b)); если $\alpha$~--- \emph{единственный} корень $f$ в $E$,
то единственный $F$-автоморфизм~--- тождественный, т.\,е.\ $\Aut(E/F)=1$.

(1) Пусть $\sqrt[3]2$~--- вещественный кубический корень из $2$. В $\Q[\sqrt[3]2][X]$
\[
X^3-2=(X-\sqrt[3]2)\,(X^2+\sqrt[3]2\,X+\sqrt[3]4),
\]
а у второго множителя нет вещественных корней (его дискриминант $-3\sqrt[3]4<0$). Так как $\Q[\sqrt[3]2]\subset\R$, единственный
корень $X^3-2$ в $\Q[\sqrt[3]2]$~--- это $\sqrt[3]2$, и $\Aut(\Q[\sqrt[3]2]/\Q)=1$, хотя $[\Q[\sqrt[3]2]:\Q]=3$. Здесь $E$ не
является полем разложения.

(2) Пусть $\operatorname{char}F=p\ne0$, $a\in F$ не является $p$-й степенью, и $E=F[\alpha]$, $\alpha^p=a$. Тогда
$X^p-a=(X-\alpha)^p$ в $E[X]$ (\rf{2.3}), так что $E$~--- поле разложения $X^p-a$, но этот многочлен имеет в $E$ лишь один
корень, и $\Aut(E/F)=1$, хотя $[E:F]=p$. Здесь $E$~--- поле разложения, но несепарабельного многочлена.
\end{example}

\begin{remark}[3.3a]
Рассуждение \rf{3.3} даёт общую формулу для простого расширения: \emph{если $E=F[\alpha]$, то $|\Aut(E/F)|$ равно числу
различных корней минимального многочлена $\mu_\alpha$ в $E$} (по \rf{2.1}(b) это число $F$-гомоморфизмов $E\to E$, а все они
автоморфизмы, так как $E$ конечномерно). Промежуточный пример: $E=\Q[\sqrt[4]2]\subset\R$. Над $E$
\[
X^4-2=(X-\sqrt[4]2)(X+\sqrt[4]2)(X^2+\sqrt2),
\]
и корни $\pm i\sqrt[4]2$ последнего множителя не вещественны. Поэтому у $\mu=X^4-2$ (неприводим по Эйзенштейну) в $E$
два корня, и $|\Aut(E/\Q)|=2$: это тождество и $\sqrt[4]2\mapsto-\sqrt[4]2$, тогда как $[E:\Q]=4$. Во всех примерах
$|\Aut(E/F)|$ делит $[E:F]$; это верно всегда (упражнение 3-5).
\end{remark}

\subsection*{Неподвижное поле и теорема Артина}

\begin{definition}[3.3b]
Пусть $G$~--- группа автоморфизмов поля $E$ (подгруппа $\Aut(E)$). Положим
\[
E^G=\operatorname{Inv}(G)=\{\alpha\in E\mid \sigma\alpha=\alpha\ \text{для всех}\ \sigma\in G\}.
\]
Это подполе $E$: если $\sigma\alpha=\alpha$ и $\sigma\beta=\beta$, то $\sigma(\alpha\pm\beta)=\alpha\pm\beta$,
$\sigma(\alpha\beta)=\alpha\beta$ и $\sigma(\alpha^{-1})=\sigma(\alpha)^{-1}=\alpha^{-1}$ при $\alpha\ne0$; кроме того, $0,1\in E^G$.
Оно называется \emph{неподвижным полем} группы $G$.
\end{definition}

То же определение имеет смысл для любого \emph{множества} $S$ автоморфизмов, и $E^S=E^{\langle S\rangle}$ (элемент,
неподвижный относительно $\sigma$ и $\tau$, неподвижен относительно $\sigma\tau$ и $\sigma^{-1}$).

\begin{remark}[3.3c: формальные свойства]
Сопоставим подгруппе $H\subset\Aut(E)$ поле $E^H$, а подполю $M\subset E$~--- группу $\Aut(E/M)$. Прямо из определений:
\begin{enumerate}[label=\textup{(\roman*)}]
\item оба сопоставления обращают включения: $H_1\subset H_2\Rightarrow E^{H_1}\supset E^{H_2}$ и
$M_1\subset M_2\Rightarrow\Aut(E/M_1)\supset\Aut(E/M_2)$;
\item $H\subset\Aut(E/E^H)$ (каждый элемент $H$ неподвижно оставляет $E^H$) и $M\subset E^{\Aut(E/M)}$ (каждый элемент
$M$ неподвижен относительно $\Aut(E/M)$);
\item поэтому $E^{\Aut(E/E^H)}=E^H$: включение $\supset$~--- это (ii) для $M=E^H$, а $\subset$ получается применением (i) к
$H\subset\Aut(E/E^H)$. Аналогично $\Aut(E/E^{\Aut(E/M)})=\Aut(E/M)$.
\end{enumerate}
Вся содержательная часть теории Галуа~--- в том, когда включения (ii) являются равенствами. Для конечной группы $H$
это так всегда (\rf{3.5}); для подполя $M$~--- ровно тогда, когда $E/M$~--- расширение Галуа (\rf{3.10}).
\end{remark}

\begin{keepfig}
\begin{center}
\begin{tikzpicture}[every node/.style={font=\small}]
\node[draw,rounded corners,cold,minimum width=4.4cm,minimum height=1.9cm] (G) at (0,0) {};
\node[above] at (G.north) {подгруппы $H\subset\Aut(E)$};
\node[draw,rounded corners,cold,minimum width=4.4cm,minimum height=1.9cm] (K) at (7.4,0) {};
\node[above] at (K.north) {подполя $M\subset E$};
\node (H) at (0,0.35) {$H$};
\node (AH) at (0,-0.45) {$\Aut(E/E^H)\supset H$};
\node (EH) at (7.4,0.35) {$E^H$};
\node (M) at (7.4,-0.45) {$M\subset E^{\Aut(E/M)}$};
\draw[->,cnew,thick] (H.east) to[bend left=10] node[above] {$H\mapsto E^H$} (EH.west);
\draw[->,cnew,thick] (6.0,-0.45) to[bend left=10] node[below] {$M\mapsto\Aut(E/M)$} (1.55,-0.45);
\end{tikzpicture}
\end{center}
\figcap{Соответствие между подгруппами и подполями. Формально верны лишь включения $H\subset\Aut(E/E^H)$ и
$M\subset E^{\Aut(E/M)}$; первое обращается в равенство для конечных $H$ (\rf{3.5}), второе~--- для $M$, над которым $E$
является расширением Галуа (\rf{3.10}).}
\end{keepfig}

\begin{theorem}[3.4: Артин]
Пусть $G$~--- конечная группа автоморфизмов поля $E$. Тогда
\[
[E:E^G]\le(G:1).
\]
\end{theorem}

\emph{Идея.} Нужно доказать, что любые $n>m=(G:1)$ элементов $E$ линейно зависимы над $F=E^G$. Линейную зависимость
над большим полем $E$ найти легко: это решение однородной системы, в которой уравнений меньше, чем неизвестных. Остаётся
<<спустить>> коэффициенты из $E$ в $F$; для этого используется то, что $G$~--- группа, и поэтому каждый $\sigma\in G$ лишь
переставляет уравнения системы.

\begin{proof}
Пусть $F=E^G$ и $G=\{\sigma_1,\dots,\sigma_m\}$, где $\sigma_1=\id$. Достаточно показать, что при $n>m$ любые элементы
$\alpha_1,\dots,\alpha_n\in E$ линейно зависимы над $F$. Рассмотрим однородную систему линейных уравнений
\begin{equation}\label{eq:artin-system}
\begin{aligned}
\sigma_1(\alpha_1)X_1+\dots+\sigma_1(\alpha_n)X_n&=0\\
&\ \,\vdots\\
\sigma_m(\alpha_1)X_1+\dots+\sigma_m(\alpha_n)X_n&=0
\end{aligned}
\end{equation}
с коэффициентами из $E$; её решения мы ищем \emph{в $E$}, а не в $F$. Уравнений $m$, неизвестных $n>m$, поэтому
пространство решений $V\subset E^n$ ненулевое (линейная алгебра над полем $E$). Выберем ненулевое решение
$(c_1,\dots,c_n)$ с наименьшим возможным числом ненулевых координат. Перенумеровав $\alpha_i$, можно считать $c_1\ne0$, а
умножив решение на $c_1^{-1}$ (решения образуют $E$-подпространство), можно считать $c_1=1\in F$. Покажем, что тогда все
$c_i\in F$. Это даст требуемое: первое уравнение системы (\ref{eq:artin-system}) (с $\sigma_1=\id$) превращается в
$\alpha_1c_1+\dots+\alpha_nc_n=0$~--- нетривиальное соотношение с коэффициентами из $F$ (ведь $c_1=1\ne0$).

Предположим, что не все $c_i$ лежат в $F$: найдутся $k$ и $i$ с $\sigma_k(c_i)\ne c_i$ (здесь $k\ne1$, так как $\sigma_1=\id$,
и $i\ne1$, так как $c_1=1$). Применим $\sigma_k$ к равенствам $\sum_j\sigma_l(\alpha_j)c_j=0$, $l=1,\dots,m$:
\[
\sum_j(\sigma_k\sigma_l)(\alpha_j)\,\sigma_k(c_j)=0,\qquad l=1,\dots,m.
\]
Так как $G$~--- группа, $\sigma_k\sigma_1,\dots,\sigma_k\sigma_m$~--- это те же $\sigma_1,\dots,\sigma_m$, только в другом порядке.
Поэтому полученные равенства~--- это уравнения системы (\ref{eq:artin-system}), записанные в другом порядке, т.\,е.\
$(\sigma_k(c_1),\dots,\sigma_k(c_n))=(1,\sigma_k(c_2),\dots,\sigma_k(c_n))$~--- тоже решение. Разность двух решений
\[
(c_1,\dots,c_n)-(1,\sigma_k(c_2),\dots,\sigma_k(c_n))=(0,\ c_2-\sigma_k(c_2),\ \dots,\ c_n-\sigma_k(c_n))
\]
--- снова решение. Оно ненулевое ($i$-я координата $c_i-\sigma_k(c_i)\ne0$), его нулевые координаты включают все нулевые
координаты $(c_j)$ (если $c_j=0$, то и $c_j-\sigma_k(c_j)=0$) и ещё первую. Значит, ненулевых координат у него меньше,
чем у $(c_j)$,~--- противоречие с выбором $(c_j)$.
\end{proof}

\begin{remark}[3.4a]
(1) Существенно, что $G$~--- группа: именно поэтому $\sigma_k$ переставляет уравнения. Для произвольного множества $S$
автоморфизмов оценка $[E:E^S]\le|S|$ неверна. Например, пусть $E=\F_8=\F_2[X]/(X^3+X+1)$ и $\varphi(x)=x^2$~---
эндоморфизм Фробениуса (автоморфизм конечного поля). Для $S=\{\id,\varphi\}$ поле $E^S$ состоит из корней $X^2-X$, т.\,е.\
$E^S=\F_2$, и $[E:E^S]=3>2=|S|$. Конечно, $E^S=E^{\langle\varphi\rangle}$, а группа $\langle\varphi\rangle$ имеет порядок $3$.

(2) Решения ищутся в $E$, а не в $F$; без этого не было бы гарантии, что система имеет ненулевое решение. Весь смысл
доказательства~--- показать, что \emph{кратчайшее} решение, нормированное условием $c_1=1$, неподвижно относительно
$G$, т.\,е.\ уже лежит в $F^n$. Иначе говоря: пространство решений $V\subset E^n$ переводится в себя покоординатным
действием каждого $\sigma\in G$; нормированный ($c_1=1$) вектор $V$ с наименьшим множеством ненулевых координат
однозначно определяется этим множеством (разность двух таких векторов с одним и тем же множеством ненулевых координат
имела бы меньше ненулевых координат), а $\sigma$ это множество не меняет. Поэтому такой вектор $G$-инвариантен.
\end{remark}

\begin{corollary}[3.5]
Пусть $G$~--- конечная группа автоморфизмов поля $E$. Тогда
\[
G=\Aut(E/E^G),
\]
и $[E:E^G]=(G:1)$.
\end{corollary}

\begin{proof}
Включение $G\subset\Aut(E/E^G)$ очевидно (\rf{3.3c}(ii)). По \rf{3.4} расширение $E/E^G$ конечно, поэтому по \rf{2.14} число
$E^G$-гомоморфизмов $E\to E$ не больше $[E:E^G]$. Получаем цепочку
\[
[E:E^G]\ \overset{\text{\rf{3.4}}}{\le}\ (G:1)\ \le\ (\Aut(E/E^G):1)\ \overset{\text{\rf{2.14}}}{\le}\ [E:E^G].
\]
Все неравенства~--- равенства. Значит, $G\subset\Aut(E/E^G)$~--- конечные группы одного порядка, т.\,е.\ они совпадают,
и $[E:E^G]=(G:1)$.
\end{proof}

\begin{remark}[3.5a: это не тавтология]
Утверждение \rf{3.5} может показаться тавтологией (<<автоморфизмы, сохраняющие неподвижное поле, \dots>>), но это не
так: $G$~--- \emph{произвольная} конечная подгруппа $\Aut(E)$, а утверждение состоит в том, что \emph{каждый}
автоморфизм $E$, тождественный на $E^G$, уже лежит в $G$. Формально (<<теоретико-множественно>>) верно лишь включение
$G\subset\Aut(E/E^G)$. Обратное включение для отдельного элемента почти очевидно: если $\varphi\in\Aut(E/E^G)$ и
$\alpha\in E$, то многочлен
\[
P_\alpha(X)=\prod_{\beta\in G\alpha}(X-\beta)
\]
(произведение по орбите $\alpha$) имеет коэффициенты в $E^G$, так как каждый $\sigma\in G$ переставляет орбиту; поэтому
$\varphi(\alpha)$~--- корень $P_\alpha$, т.\,е.\ $\varphi(\alpha)=\sigma(\alpha)$ для некоторого $\sigma\in G$. Трудность в том,
что $\sigma$ зависит от $\alpha$, а нужен один $\sigma$, годный сразу для всех $\alpha$. (Если бы $E$ было порождено над
$E^G$ одним элементом $\gamma$, хватило бы взять $\sigma$ с $\sigma\gamma=\varphi\gamma$.) Подсчёт степеней в доказательстве
\rf{3.5} обходит эту трудность. Что конечность $G$ здесь существенна, показывает следующий пример.
\end{remark}

\begin{example}[3.5b]
Пусть $\operatorname{char}k=0$, $E=k(X)$ и $G=\langle\sigma\rangle$, где $\sigma(f)(X)=f(X+1)$; группа $G$ бесконечна (так
как $\sigma^n(X)=X+n\ne X$ при $n\ne0$). \emph{Тогда $E^G=k$.} Пусть $r=g/h\in E^G$, где $g,h\in k[X]$ взаимно просты и $h$
унитарен. Из $g(X+1)h(X)=g(X)h(X+1)$ следует, что $h(X)$ делит $g(X)h(X+1)$, а значит, и $h(X+1)$; оба многочлена унитарны
одной степени, поэтому $h(X+1)=h(X)$ и затем $g(X+1)=g(X)$. Но если многочлен $q$ удовлетворяет $q(X+1)=q(X)$, то
$q(X)-q(0)$ обращается в нуль в точках $0,1,2,\dots$, которые в характеристике $0$ попарно различны; значит, $q$~---
константа. Итак, $g,h\in k$ и $r\in k$.

Следовательно, $[E:E^G]=\infty$ (оценка \rf{3.4} невозможна), а $\Aut(E/E^G)=\Aut(k(X)/k)\cong PGL_2(k)$ (\rf{3.1b})
гораздо больше $G$: например, $X\mapsto 2X$ не лежит в $G$. Для бесконечных групп правильная формулировка требует
топологии (глава 7).

В характеристике $p$ картина иная: $\sigma^p(X)=X+p=X$, группа $G$ имеет порядок $p$, и $E^G=k(X^p-X)$ (\rf{3.5c}(2)).
\end{example}

Теорема Артина даёт удобный способ находить неподвижные поля: \emph{если $u_1,\dots,u_r\in E^G$ и
$[E:k(u_1,\dots,u_r)]\le(G:1)$ для подполя $k\subset E^G$, то $E^G=k(u_1,\dots,u_r)$}. Действительно,
$k(u_1,\dots,u_r)\subset E^G$, а по \rf{3.5} $[E:E^G]=(G:1)\ge[E:k(u_1,\dots,u_r)]$, так что поля совпадают (\rf{1.20}).

\begin{example}[3.5c]
Пусть $k$~--- поле, $E=k(X)$; степени расширений вычисляем по \rf{3.1a}.
\begin{enumerate}[label=\textup{(\arabic*)}]
\item $G=\{\id,\ X\mapsto1/X\}$ (порядок $2$, так как $1/X\ne X$). Элемент $u=X+1/X$ неподвижен, а
$u=\frac{X^2+1}{X}$ имеет степень $2$. Значит, $k(X)^G=k(X+1/X)$.
\item $\operatorname{char}k=p$, $G=\langle X\mapsto X+1\rangle$ порядка $p$. Элемент $u=X^p-X$ неподвижен:
$(X+1)^p-(X+1)=X^p-X$. Его степень $p$, поэтому $k(X)^G=k(X^p-X)$. Здесь $X$~--- корень многочлена
Артина--Шрайера $T^p-T-u$ над $k(u)$ (ср.\ \rf{2.22b}).
\item Шесть подстановок
\[
X,\quad \frac1X,\quad 1-X,\quad \frac1{1-X},\quad \frac{X-1}X,\quad \frac X{X-1}
\]
образуют группу $G\cong S_3$ (она переставляет точки $0,1,\infty$ прямой; композиции этих подстановок снова среди них,
а соответствующие матрицы попарно непропорциональны над любым полем). Функция
\[
j=\frac{(X^2-X+1)^3}{X^2(X-1)^2}
\]
неподвижна относительно образующих $X\mapsto1/X$ и $X\mapsto1-X$, а значит, и относительно $G$; её степень $6=(G:1)$.
Поэтому $k(X)^G=k(j)$.
\end{enumerate}
\end{example}
